% A.91 reserved by the root.
\subsection{Contact four fixes the dyadic support without a pairing assumption}

\begin{statement}{A.91}{Proved}
{The full contact-four comparison and the family $\pm16p^n$}
In any labeled integral realization of A.14 whose rational
parameter has dyadic contact order four, the actual ordinary
dyadic point has exactly the same valuation pair as that
parameter.  No condition on the global Weil multiplier
is required.

Consequently, within their original parameter ranges, the criteria of A.79 and A.82 apply with an
arbitrary labeled integral comparison whenever
$v_2(\lambda)=4$; their full-$H$ proviso at this valuation
can be omitted.
For every odd prime $p$, every integer $n\ge1$ and every
$\epsilon\in\{1,-1\}$, the parameter
\[
                             \lambda=\epsilon16p^n
\tag{A.91.1}
\]
is recognized with an arbitrary labeled integral realization.
The same holds for all six fractional-linear transforms
permuting $0,1,\infty$.  This removes the geometric-conjugator
requirement from the family in A.64.
\end{statement}
\begin{proof}
For $v\in\Z_2^\times$, denote by $G_v$ the good curve
\[
                G_v:\quad Y^2+XY=X^3+4vX^2+vX.
\tag{A.91.2}
\]
The exact change of variables (A.87.8) identifies
$G_v$ with $E_{16v}^{(-1)}$, and its discriminant
$v^2(1-16v)^2$ is a unit.
For every odd coefficient $\ell$, its inertia action on
$T_\ell$ is therefore the identity.
Let $\chi_d$ denote the quadratic character of
$\Q_2(\sqrt d)/\Q_2$.
The explicit twists (A.87.3)--(A.87.4) give the following
table, where the squareclasses in the middle column have
their powers of sixteen removed:
\[
\begin{array}{c|c|c}
 t & ([t],[1-t])\text{ in }(\Q_2^\times/(\Q_2^\times)^2)^2
                              & T_\ell(E_t)|_{I_2}\\ \hline
 16v &(v,1)&\chi_{-1}\cdot I\\
 1-16v &(1,v)&I\\
 v/16 &(v,-v)&\chi_{-v}\cdot I .
\end{array}
\tag{A.91.3}
\]
For the last row, $E_{v/16}=E_{16/v}^{(v/16)}$ is the
$(-v/16)$-twist of the good curve
$E_{16/v}^{(-1)}=G_{1/v}$; the factor sixteen is a
square, giving precisely $\chi_{-v}$.
For the middle row, the translation
$(x,y)\mapsto(x-1,y)$ identifies $E_{1-16v}$ with
$E_{16v}^{(-1)}$ itself.
The character $\chi_{-1}$ is nontrivial on inertia:
in $\Q_2(i)$ the element $1+i$ has valuation $1/2$
because $(1+i)^2=2i$, so this quadratic extension
is ramified.  Thus the first row has inertia order
exactly two, while the middle row has order one.

The actual local section is ordinary by A.37 and its
point curve is potentially good.  A.16 supplies the
same two dyadic Kummer squareclasses as the realized
parameter.  Their even valuation parities leave only
contact orders two and four.
A.87 gives inertia order twenty-four at contact two,
whereas every row in (A.91.3) has inertia order at
most two.  The integral intertwiner excludes contact
two for the actual point.  It remains to check the
three contact-four rows with the precise squareclasses.

Suppose first that the realized parameter is $16u$.
An actual contact-one point has unit classes $(1,v)$.
Equality with $(u,1)$ forces $[u]=[v]=1$, but its
inertia is trivial rather than $\chi_{-1}$.
An actual contact-infinity point has classes $(v,-v)$.
Equality with $(u,1)$ forces $[-v]=1$.
Its inertia is then also trivial, by the last row of
(A.91.3), again impossible.
Only contact zero remains, proving the exact pair
$(4,0)$ for every odd unit $u$, including negative
rational units.

For a realized contact-one parameter $1-16u$, an
actual contact-zero point would have nontrivial
$\chi_{-1}$ inertia, unlike the realized curve.
An actual contact-infinity point can have the same
classes $(1,u)$ only when $[v]=1$ and $[u]=[-1]$.
Its inertia is then $\chi_{-v}=\chi_{-1}$ on inertia,
also impossible.  Thus the pair is $(0,4)$.
Finally take a realized contact-infinity parameter
$u/16$.  An actual contact-zero point would require
$[-u]=1$, so the realized inertia would be trivial
while that actual point has inertia $\chi_{-1}$.
An actual contact-one point would require $[u]=1$,
so the realized inertia would be $\chi_{-1}$ while
that actual point is good.  Both are excluded, and
the pair is $(-4,-4)$.
These comparisons prove the first assertion at all
labels using the actual scalar inertia characters.

For an even integer parameter with $v_2(\lambda)=4$,
the result supplies $v_2(A)=4$ and $v_2(B)=0$.
This is exactly the support information furnished by
the full-$H$ assumption in A.79 and A.82 at $k=4$.
Every later step of their proofs retains an arbitrary
integral multiplier, so the stated extension follows.

We give the resulting recognition proof for (A.91.1)
to identify precisely where its earlier assumption
has disappeared.  The dyadic result and the ordered
Tate-line calculation at all odd bad primes give
\[
                         A=\epsilon16p^\beta,
                              \qquad\beta\in\widehat\Z.
\tag{A.91.4}
\]
Indeed, no odd prime other than $p$ contributes to
the first coordinate; at such a bad prime the label
is one.  The sign in (A.91.4) is fixed by the global
Kummer squareclass in A.16.
At $p$ a possible cusp has label zero.  Its exact
formal Tate extension class is
\[
                              A^2/256=p^{2\beta}.
\tag{A.91.5}
\]
At coefficient $p$ it lies on the rational Kummer
line $\Q_p\kappa_p(p)$.
The rational-line transcendence theorem A.63
excludes the realized ordinary Tate class from
this line, even after a common finite splitting
extension and under an arbitrary nonzero scalar
comparison.  Thus this local section cannot be a
cusp.  It is ordinary, and A.78 gives its exact
parameter $\lambda$.
The full local Kummer valuation in (A.91.4) then
gives $\beta=n$ in $\widehat\Z$, not merely in
one coefficient.  Consequently $A=\lambda$ is
one rational number.  The finite-value recognition
argument of A.64, or equivalently A.14 and A.37,
now makes the full section the point section
$s_\lambda$.

We finish by spelling out the transfer to all six
parameters.  The two generators of their permutation
group are $\sigma(t)=1-t$ and $\iota(t)=1/t$.
The maps on the actual Legendre curves are
\[
 (x,y)\longmapsto(1-x,iy)
          \quad\text{from }E_t\text{ to }E_{1-t},
 \qquad
 (x,y)\longmapsto(x/t,y/(t\sqrt t))
          \quad\text{from }E_t\text{ to }E_{1/t}.
\tag{A.91.6}
\]
They are defined over the indicated quadratic
extensions and preserve the Weil pairing as curve
isomorphisms.  Their Galois descent factors are
respectively $\chi_{-1}$ and $\chi_t$.
They permute the ordered nonzero two-torsion by
the same permutation on the section and on its
realized point.  The first factor is constant;
the second agrees on these two arithmetic sections
because A.16 gives $[A]=[\lambda]$ modulo squares.
Thus the existing integral intertwiner, transported
through these actual curve maps, still intertwines
the same quadratic twist on both sides.
Canceling that common character gives a labeled
integral realization for the transformed section.
No geometric-conjugator hypothesis enters this
construction.  Applying it to the inverse
permutation reduces every transformed parameter
to (A.91.1); the inverse automorphism of $U$ then
returns exactly its own rational point section.
\end{proof}

The symbolic good model, the full table of dyadic
unit squareclasses and the negative-unit cases are
checked in
\path{code/legendre_dyadic_four_certificate.py},
with output in
\path{code/legendre_dyadic_four_certificate_output.txt}.
